% !TeX root = main.tex
\chapter*{Reading alongside Lorentzian Sewing}\addcontentsline{toc}{chapter}{Reading alongside Lorentzian Sewing}\markboth{Companion guide}{}
The bold numbers below refer to the chapters or appendices of \emph{Lorentzian Sewing: Geometry, Dynamics, and Observables}, the volume this text accompanies. Chapter, section, equation and page references after them refer to the present book. This distinction matters: a theorem giving the conditions for a construction is not the assertion that those conditions hold in every proposed field theory.

\begin{description}[style=nextline,leftmargin=0pt,labelindent=0pt,labelwidth=0pt,labelsep=0pt,font=\normalfont\bfseries]
\item[1. Sewing spacetime regions.] For the quotient topology use Theorem~\ref{f:q-descent}; for an actual smooth atlas use the sector and overlap test on p.~\pageref{f:sector-theorem}. Regular levels, normal quotients and surviving exterior walls are built in Chapter~\ref{f:submanifolds}. A causal test is separate: apply the shifted cylinder theorem on p.~\pageref{f:shift-cylinder}.
\item[2. Bundles and connections at a cut.] The bundle-gluing theorem on p.~\pageref{f:bundle-gluing} constructs fibers, not merely their coordinate values. Check a proposed lift by \eqref{f:morphism-overlap}; compare connection forms only after pulling to the same base, as in \eqref{f:connection-change}. The two sign orbits and their stabilizers appear beside \eqref{f:sign-groupoid}.
\item[3. Smooth fields and gauge transformations.] Start with the affine space of first jets in \S\ref{f:jet}. Proposition~\ref{f:jet-pasting} gives necessity and sufficiency for \emph{actual} half-sections, including mixed orders and the condition at a remaining exterior wall. The same test is applied to the arrows of the gauge groupoid, not to a set from which the arrows have already been removed.
\item[4. Varying a region.] Localize and integrate an actual top form using Theorem~\ref{f:partition}; Stokes is proved on p.~\pageref{f:stokes}. In \S\ref{f:first-variation} the action differentiates on vertical field variations and yields both an Euler operator and a boundary covector. Fixed and released trace problems therefore have different stationarity conditions.
\item[5. Covariant phase space and boundary charges.] Chapter~\ref{f:variation} keeps the wall action and its corner potential in the integrated slice potential. Chapter~\ref{f:symplectic} derives the Hamiltonian convention and the Noether charge, and Theorem~\ref{f:regular-reduction} states the regular hypotheses for reduction. The gauge-complex boundary exercise on p.~\pageref{f:gauge-boundary-ex} tests whether a parameter is proper.
\item[6. Causal response and the Peierls bracket.] Pair the Hessian with dual densities by \eqref{f:green-variation}; check both support and adjoint domains in \S\ref{f:green-domains}, including the one-sided extension on p.~\pageref{f:one-sided-extension}. Compact source representatives and their Hamiltonian solutions are related by Theorem~\ref{f:source-exact}. For a gauge equation the relevant inverse is instead the degree-lowering homotopy of Proposition~\ref{f:green-homotopy}.
\item[7. Transmission and reconstruction.] Chapter~\ref{f:variation} recovers normal derivatives from value and conjugate momentum on one inverse branch, and gives a counterexample with two regular branches. Compare that local statement with the symbol and the evolution estimates in \S\ref{f:wave-symbol} and \S\ref{f:wave-energy}. The last exercise of Chapter~\ref{f:variation} distinguishes choosing a gauge from transmitting an observable derivative.
\item[8. Sewing a causal response.] The trace and its right inverse on a flat face are in Theorem~\ref{f:plane-trace}, and the weak flux is in Proposition~\ref{f:weak-flux}. The joined inverse on p.~\pageref{f:sewn-inverse} is conditional on an inverse on the produced mismatches and on causal uniqueness. The impedance example realizes the required history map rather than assuming invertibility of an abstract matrix.
\item[9. Presymplectic sewing and observables.] A kernel is taken \emph{after} restricting the form (Example~\ref{f:two-rotor}). In Chapter~\ref{f:quantum} the invariant algebra of two open transports gives a closed loop observable; preserving brackets among already known generators need not generate this observable.
\item[10. Junctions, elimination and refinement.] Use \S\ref{f:glue-space} for simultaneous sectors, then \S\ref{f:delayed-schur} for Schur elimination on histories. The latter keeps the unsolved boundary data as variables. Geometric compatibility alone says nothing about whether an inverse pivot preserves the direction of travel in time.
\item[11. The linear commutator.] Apply \eqref{f:label-sign} on source labels, not on unreduced field components. Chapter~\ref{f:quantum} explicitly constructs a CCR algebra, a state and its GNS representation as three operations. Proposition~\ref{f:gaussian-positivity} tests candidate mixed covariances; Corollary~\ref{f:wick-pairings} derives all ordered moments of a centered quasifree state. The operator-algebra meaning of a sharp cut is discussed on p.~\pageref{f:local-factors}.
\item[12. Perturbations about a sewn solution.] Linearization and domains come from Chapters~\ref{f:variation} and \ref{f:waves}. Equation~\eqref{f:ext-retarded} separates the background, the free inverse and the nonlinear source. Equation~\eqref{f:ext-tree} then calculates a particular mixed response; the next section compares it with ordered insertions. Formal orders are algebraic coefficients, not a statement of convergence at finite coupling.
\item[13. Time-ordered products.] Theorems~\ref{f:wf-pullback} and \ref{f:wf-hadamard} specify singular directions and permissible pullbacks; \eqref{f:ext-fish} is a first product off the diagonal. Theorem~\ref{f:ext-theorem} bounds the freedom at coincidence. The scale residue \eqref{f:ext-residue} and Ward contact \eqref{f:ext-ward} are distinct constraints; neither assigns physical weights to global sectors.
\end{description}

\paragraph{The earlier appendices.} For A (orientations and flux), use Chapters~\ref{f:forms} and \ref{f:variation}. For B (traces and hyperbolic domains), use Chapters~\ref{f:distributions}--\ref{f:waves}. For C (Hadamard tests), the geodesic input is constructed in Chapter~\ref{f:connections} before the theorem on p.~\pageref{f:wf-hadamard}; the singular-product and extension arguments are in Chapters~\ref{f:microlocal} and \ref{f:extension}. The two-strip construction of D reappears in Chapter~\ref{f:capstone}, where it determines three different kernels. Compare the static response \eqref{f:static-measurement} with the difference of subtracted vacuum fluctuations \eqref{f:relative-fluctuation}.

\paragraph{What remains an input.} Among the results for which we supply hypotheses and an application rather than a full general proof are the smooth collar and convex-neighborhood existence theorems, the global spacetime splitting and normally hyperbolic Cauchy theorems, Plancherel's theorem, the free proper action quotient theorem, and the microlocal pullback, propagation, Hadamard and extension theorems. The text does prove, for example, the flat sharp trace bound, finite-energy uniqueness in the smooth flat model, the compact-source exactness argument, GNS, and the local counterterm ambiguity once the corresponding inputs are available. No general boundary Green theorem, proof of quantization of interacting gauge fields, or classification of physical defects is assumed to follow from those results.

% !TeX root = main.tex
\chapter*{Index of concepts}\addcontentsline{toc}{chapter}{Index of concepts}\markboth{Index of concepts}{}
The page given is a place to learn or use the concept, not merely its first occurrence. The PDF page links lead to the corresponding construction.
\newcommand{\fentry}[2]{\par\noindent\hangindent=1em\hangafter=1 #1, \pageref{#2}\par}
\newcommand{\fcontinued}[3]{\par\noindent\hangindent=1em\hangafter=1 #1, \pageref{#2}, \pageref{#3}\par}
\begingroup\setlength{\columnsep}{1.7em}\raggedright
\subsection*{A--G}
\begin{multicols}{2}
\fentry{Adjoint boundary domain}{f:green-domains}
\fentry{Affine space of derivatives}{f:jet}
\fentry{Algebra by generators and relations}{f:polynomial-ccr}
\fentry{Associated bundle}{f:associated-bundle}
\fentry{Atlas and chart changes}{f:stereo}
\fentry{Basis of a topology}{f:top-basis}
\fcontinued{Boundary action and charge}{f:first-variation}{f:gauge-boundary-ex}
\fentry{Boundary value, distributional}{f:wf-bv}
\fentry{Bundle map over a base map}{f:morphism-overlap}
\fentry{Bump function}{f:partition}
\fentry{Canonical commutation relations}{f:polynomial-ccr}
\fentry{Causal factorization}{f:ext-factor}
\fentry{Causal propagator}{f:causal-propagator}
\fentry{Cauchy temporal function}{f:cauchy-splitting}
\fentry{Characteristic hypersurface}{f:wave-symbol}
\fentry{Clutching and cocycles}{f:bundle-gluing}
\fentry{Cohomology and chain homotopy}{f:gauge-complex}
\fentry{Collar and corner}{f:glue-space}
\fentry{Compactness and separation}{f:q-compact-hausdorff}
\fentry{Completion and Hilbert dual}{f:hilbert-representation}
\fentry{Connected component}{f:q-components}
\fentry{Connection, change of frame}{f:connection-change}
\fcontinued{Conormal space, conormal test}{f:normal}{f:wf-pullback}
\fentry{Constraint, regular reduction}{f:regular-reduction}
\fentry{Continuity and homeomorphism}{f:top-continuity}
\fentry{Convex normal neighborhood}{f:convex-normal}
\fentry{Counterterms}{f:ext-jets}
\fentry{Covariance, positivity}{f:gaussian-positivity}
\fentry{Curvature}{f:curvature}
\fentry{Deformation complex}{f:deformation-complex}
\fentry{Density and orientation}{f:orientation-density}
\fentry{Derivation at a point}{f:tangent-coords}
\fentry{Differential forms}{f:forms}
\fentry{Distribution, order and support}{f:distribution-estimate}
\fentry{Energy and finite propagation}{f:energy-estimate}
\fentry{Euler expression}{f:first-variation}
\fentry{Exact sequence of source labels}{f:source-exact}
\fentry{Exponential map}{f:convex-normal}
\fentry{Extension across a collision}{f:ext-theorem}
\fentry{Exterior derivative and Lie derivative}{f:exterior}
\fentry{Factor and split inclusion}{f:local-factors}
\fentry{Fej\'er kernel}{f:circle-fourier}
\fentry{Feynman contraction}{f:ext-feynman}
\fentry{Fluctuation, subtracted difference}{f:relative-fluctuation}
\fentry{Fourier transform, convention}{f:fourier-pair}
\fentry{Gauge fixing}{f:green-homotopy}
\fentry{Gauss flux at a boundary}{f:gauge-boundary-ex}
\fentry{Geodesic and world function}{f:world-function}
\fentry{Germs and tangents}{f:tangent-coords}
\fentry{Global hyperbolicity}{f:cauchy-splitting}
\fentry{GNS representation}{f:gns}
\fentry{Green operators, compact and one-sided domains}{f:green-identities}
\fentry{Groupoid and stabilizer}{f:sign-groupoid}
\end{multicols}\clearpage
\subsection*{H--W}
\begin{multicols}{2}
\fcontinued{Hadamard state and parametrix}{f:wf-hadamard}{f:circle-subtraction}
\fentry{Hamiltonian vector and bracket sign}{f:poisson-convention}
\fentry{Hausdorff failure after identification}{f:q-double-origin}
\fentry{Hodge star and formal adjoint}{f:hodge}
\fentry{Holonomy}{f:parallel-ode}
\fentry{Hypersurface, regular value}{f:regular-value}
\fentry{Jet bundle and holonomic jet}{f:jet}
\fentry{Levi--Civita connection}{f:levi-civita}
\fentry{Lie algebra and representation}{f:group-action}
\fentry{Lindel\"of property and second countability}{f:q-lindelof}
\fentry{Momentum map}{f:regular-reduction}
\fentry{Normal Legendre branch}{f:normal-legendre}
\fentry{Null screen}{f:null-screen}
\fentry{Open set and neighborhood}{f:top-open}
\fentry{Orbit, stabilizer, transition signs}{f:sign-groupoid}
\fentry{Parallel transport}{f:parallel-ode}
\fentry{Partition of unity}{f:partition}
\fentry{Peierls bracket}{f:peierls-labels}
\fentry{Presymplectic kernel}{f:symplectic}
\fentry{Principal symbol}{f:wave-symbol}
\fentry{Propagation of singularities}{f:wf-propagation}
\fentry{Pullback of a distribution}{f:wf-pullback}
\fentry{Quasifree state}{f:gaussian-positivity}
\fentry{Quotient topology and descent}{f:q-descent}
\fentry{Radical and physical observables}{f:alternating-quotient}
\fentry{Scaling degree and ambiguity}{f:ext-theorem}
\fentry{Schur complement on histories}{f:delayed-schur}
\fentry{Sobolev spaces and duality}{f:sobolev-scale}
\fentry{Static response of a twisted field}{f:static-measurement}
\fentry{Stokes' theorem}{f:stokes}
\fentry{Tensor product and complexification}{f:tensor-product}
\fentry{Test functions}{f:test-topology}
\fcontinued{Trace, strong and weak}{f:plane-trace}{f:weak-flux}
\fentry{Transversality and clean intersection}{f:normal}
\fentry{Tree response}{f:ext-tree}
\fentry{Two-point function, mode construction}{f:circle-two-point}
\fentry{Van Vleck determinant}{f:van-vleck}
\fentry{Ward contact and consistency}{f:ext-ward}
\fentry{Wavefront set}{f:wf-test}
\fentry{Weyl algebra}{f:weyl-relations}
\fcontinued{Wick pairings and a two-line loop}{f:wick-pairings}{f:ext-fish}
\fentry{Wilson loop}{f:two-rotor}
\end{multicols}\endgroup

